\documentclass[12pt]{ximera}
\typeout{Start loading xmPreamble.tex}%

% Add here extra macro's that are loaded automatically by all documents of claas 'ximera' or 'xourse' in this repo

%%
%%  Example:
%%
% \newcommand{\R}{\mathbb{R}}
\usepackage{tikz}
\usetikzlibrary{positioning, arrows.meta}
\usepackage{multicol}
\usepackage{enumitem}
\usepackage{caption}
\usepackage{amsmath}
\usepackage{amsthm}
\usepackage{fontenc}

%% ---------------------------------------------------------------
%% Shared worksheet macros.
%%
%% These came from the Summer 2026 worksheet set, where each worksheet carried its
%% own copy. They have to live here instead: a xourse inlines every activity into a
%% single document, so duplicate \newcommand definitions across activities collide,
%% and the xourse gobbles \input so a per-topic macro file never loads.
%%
%% They use \providecommand, NOT \newcommand. ximera.cls loads this file automatically,
%% and every activity also carries an explicit \typeout{Start loading xmPreamble.tex}%

% Add here extra macro's that are loaded automatically by all documents of claas 'ximera' or 'xourse' in this repo

%%
%%  Example:
%%
% \newcommand{\R}{\mathbb{R}}
\usepackage{tikz}
\usetikzlibrary{positioning, arrows.meta}
\usepackage{multicol}
\usepackage{enumitem}
\usepackage{caption}
\usepackage{amsmath}
\usepackage{amsthm}
\usepackage{fontenc}

%% ---------------------------------------------------------------
%% Shared worksheet macros.
%%
%% These came from the Summer 2026 worksheet set, where each worksheet carried its
%% own copy. They have to live here instead: a xourse inlines every activity into a
%% single document, so duplicate \newcommand definitions across activities collide,
%% and the xourse gobbles \input so a per-topic macro file never loads.
%%
%% They use \providecommand, NOT \newcommand. ximera.cls loads this file automatically,
%% and every activity also carries an explicit \typeout{Start loading xmPreamble.tex}%

% Add here extra macro's that are loaded automatically by all documents of claas 'ximera' or 'xourse' in this repo

%%
%%  Example:
%%
% \newcommand{\R}{\mathbb{R}}
\usepackage{tikz}
\usetikzlibrary{positioning, arrows.meta}
\usepackage{multicol}
\usepackage{enumitem}
\usepackage{caption}
\usepackage{amsmath}
\usepackage{amsthm}
\usepackage{fontenc}

%% ---------------------------------------------------------------
%% Shared worksheet macros.
%%
%% These came from the Summer 2026 worksheet set, where each worksheet carried its
%% own copy. They have to live here instead: a xourse inlines every activity into a
%% single document, so duplicate \newcommand definitions across activities collide,
%% and the xourse gobbles \input so a per-topic macro file never loads.
%%
%% They use \providecommand, NOT \newcommand. ximera.cls loads this file automatically,
%% and every activity also carries an explicit \input{../xmPreamble}, so this file is
%% read twice. That was harmless while it held only \usepackage lines (LaTeX skips a
%% package it has already loaded) but a second \newcommand is a hard error.
%%
%%   \blank[width]        fill-in-the-blank rule (default 2.5cm)
%%   \showwork            "show and explain your work" reminder
%%   \showworkline        ... plus which schedule line was used
%%   \showworkyear        ... plus which year's schedule was used
%%   \schedhead{...}      centred bold caption above a tiered-rate table
%%   \samesquares[scale]{left}{right}   two equal-area squares, labelled
%%   \samecubes[scale]{left}{right}     two equal-volume cubes, labelled
%% ---------------------------------------------------------------

\providecommand{\blank}[1][2.5cm]{\underline{\hspace{#1}}}
\providecommand{\showwork}{\textbf{REMEMBER TO SHOW AND EXPLAIN YOUR WORK.}}
\providecommand{\showworkline}{\textbf{REMEMBER TO SHOW AND EXPLAIN YOUR WORK.
  Explanation should include how and why you know which line of the
  schedule was used to calculate this amount.}}
\providecommand{\showworkyear}{\begin{sloppypar}\textbf{REMEMBER TO SHOW AND
  EXPLAIN YOUR WORK. Explanation should include how and why you know which
  line of the year's tax schedule was used to calculate this tax
  amount.}\end{sloppypar}}
\providecommand{\schedhead}[1]{\begin{center}\textbf{#1}\end{center}}

\providecommand{\samesquares}[3][1]{%
\begin{center}
{\footnotesize These squares are the \textbf{same size}.}

\vspace{1.5mm}
\begin{tikzpicture}[scale=#1,every node/.style={scale=#1}]
  \draw (0,0) rectangle (2.4,2.4);
  \node[anchor=north west] at (0.15,2.25) {Area:};
  \node[anchor=east]  at (-0.15,1.2) {#2};
  \node[anchor=north] at (1.2,-0.15) {#2};
  \begin{scope}[xshift=4.6cm]
    \draw (0,0) rectangle (2.4,2.4);
    \node[anchor=north west] at (0.15,2.25) {Area:};
    \node[anchor=east]  at (-0.15,1.2) {#3};
    \node[anchor=north] at (1.2,-0.15) {#3};
  \end{scope}
\end{tikzpicture}
\end{center}}

\providecommand{\samecubes}[3][1]{%
\begin{center}
{\footnotesize These cubes are the \textbf{same size}.}

\vspace{1.5mm}
\begin{tikzpicture}[scale=#1,every node/.style={scale=#1}]
  \foreach \dx in {0,5.2} {
    \begin{scope}[xshift=\dx cm]
      \draw (0,0) rectangle (2.0,2.0);
      \draw (0,2.0) -- (0.65,2.6) -- (2.65,2.6) -- (2.0,2.0);
      \fill[black!12] (2.0,0) -- (2.65,0.6) -- (2.65,2.6) -- (2.0,2.0) -- cycle;
      \draw (2.0,0) -- (2.65,0.6) -- (2.65,2.6) -- (2.0,2.0);
      \node[anchor=north west] at (0.1,1.9) {Volume:};
    \end{scope}
  }
  \node[anchor=east]  at (-0.15,1.0) {#2};
  \node[anchor=north] at (1.0,-0.15) {#2};
  \node[anchor=west]  at (2.25,0.4) {#2};
  \node[anchor=east]  at (5.05,1.0) {#3};
  \node[anchor=north] at (6.2,-0.15) {#3};
  \node[anchor=west]  at (7.45,0.4) {#3};
\end{tikzpicture}
\end{center}}

%% NOTE: do NOT add \usepackage to individual activity .tex files.
%% When a xourse inlines an activity it gobbles \documentclass and \input, but a bare
%% \usepackage still executes -- after \begin{document} -- and errors out.
%% All shared packages and macros belong here.
%%
%% ximera.cls already loads: amsmath, amssymb, amsthm, cancel, enumitem, graphicx,
%% hyperref, listings, pgfplots, tikz, url, xcolor. No need to re-load those.
, so this file is
%% read twice. That was harmless while it held only \usepackage lines (LaTeX skips a
%% package it has already loaded) but a second \newcommand is a hard error.
%%
%%   \blank[width]        fill-in-the-blank rule (default 2.5cm)
%%   \showwork            "show and explain your work" reminder
%%   \showworkline        ... plus which schedule line was used
%%   \showworkyear        ... plus which year's schedule was used
%%   \schedhead{...}      centred bold caption above a tiered-rate table
%%   \samesquares[scale]{left}{right}   two equal-area squares, labelled
%%   \samecubes[scale]{left}{right}     two equal-volume cubes, labelled
%% ---------------------------------------------------------------

\providecommand{\blank}[1][2.5cm]{\underline{\hspace{#1}}}
\providecommand{\showwork}{\textbf{REMEMBER TO SHOW AND EXPLAIN YOUR WORK.}}
\providecommand{\showworkline}{\textbf{REMEMBER TO SHOW AND EXPLAIN YOUR WORK.
  Explanation should include how and why you know which line of the
  schedule was used to calculate this amount.}}
\providecommand{\showworkyear}{\begin{sloppypar}\textbf{REMEMBER TO SHOW AND
  EXPLAIN YOUR WORK. Explanation should include how and why you know which
  line of the year's tax schedule was used to calculate this tax
  amount.}\end{sloppypar}}
\providecommand{\schedhead}[1]{\begin{center}\textbf{#1}\end{center}}

\providecommand{\samesquares}[3][1]{%
\begin{center}
{\footnotesize These squares are the \textbf{same size}.}

\vspace{1.5mm}
\begin{tikzpicture}[scale=#1,every node/.style={scale=#1}]
  \draw (0,0) rectangle (2.4,2.4);
  \node[anchor=north west] at (0.15,2.25) {Area:};
  \node[anchor=east]  at (-0.15,1.2) {#2};
  \node[anchor=north] at (1.2,-0.15) {#2};
  \begin{scope}[xshift=4.6cm]
    \draw (0,0) rectangle (2.4,2.4);
    \node[anchor=north west] at (0.15,2.25) {Area:};
    \node[anchor=east]  at (-0.15,1.2) {#3};
    \node[anchor=north] at (1.2,-0.15) {#3};
  \end{scope}
\end{tikzpicture}
\end{center}}

\providecommand{\samecubes}[3][1]{%
\begin{center}
{\footnotesize These cubes are the \textbf{same size}.}

\vspace{1.5mm}
\begin{tikzpicture}[scale=#1,every node/.style={scale=#1}]
  \foreach \dx in {0,5.2} {
    \begin{scope}[xshift=\dx cm]
      \draw (0,0) rectangle (2.0,2.0);
      \draw (0,2.0) -- (0.65,2.6) -- (2.65,2.6) -- (2.0,2.0);
      \fill[black!12] (2.0,0) -- (2.65,0.6) -- (2.65,2.6) -- (2.0,2.0) -- cycle;
      \draw (2.0,0) -- (2.65,0.6) -- (2.65,2.6) -- (2.0,2.0);
      \node[anchor=north west] at (0.1,1.9) {Volume:};
    \end{scope}
  }
  \node[anchor=east]  at (-0.15,1.0) {#2};
  \node[anchor=north] at (1.0,-0.15) {#2};
  \node[anchor=west]  at (2.25,0.4) {#2};
  \node[anchor=east]  at (5.05,1.0) {#3};
  \node[anchor=north] at (6.2,-0.15) {#3};
  \node[anchor=west]  at (7.45,0.4) {#3};
\end{tikzpicture}
\end{center}}

%% NOTE: do NOT add \usepackage to individual activity .tex files.
%% When a xourse inlines an activity it gobbles \documentclass and \input, but a bare
%% \usepackage still executes -- after \begin{document} -- and errors out.
%% All shared packages and macros belong here.
%%
%% ximera.cls already loads: amsmath, amssymb, amsthm, cancel, enumitem, graphicx,
%% hyperref, listings, pgfplots, tikz, url, xcolor. No need to re-load those.
, so this file is
%% read twice. That was harmless while it held only \usepackage lines (LaTeX skips a
%% package it has already loaded) but a second \newcommand is a hard error.
%%
%%   \blank[width]        fill-in-the-blank rule (default 2.5cm)
%%   \showwork            "show and explain your work" reminder
%%   \showworkline        ... plus which schedule line was used
%%   \showworkyear        ... plus which year's schedule was used
%%   \schedhead{...}      centred bold caption above a tiered-rate table
%%   \samesquares[scale]{left}{right}   two equal-area squares, labelled
%%   \samecubes[scale]{left}{right}     two equal-volume cubes, labelled
%% ---------------------------------------------------------------

\providecommand{\blank}[1][2.5cm]{\underline{\hspace{#1}}}
\providecommand{\showwork}{\textbf{REMEMBER TO SHOW AND EXPLAIN YOUR WORK.}}
\providecommand{\showworkline}{\textbf{REMEMBER TO SHOW AND EXPLAIN YOUR WORK.
  Explanation should include how and why you know which line of the
  schedule was used to calculate this amount.}}
\providecommand{\showworkyear}{\begin{sloppypar}\textbf{REMEMBER TO SHOW AND
  EXPLAIN YOUR WORK. Explanation should include how and why you know which
  line of the year's tax schedule was used to calculate this tax
  amount.}\end{sloppypar}}
\providecommand{\schedhead}[1]{\begin{center}\textbf{#1}\end{center}}

\providecommand{\samesquares}[3][1]{%
\begin{center}
{\footnotesize These squares are the \textbf{same size}.}

\vspace{1.5mm}
\begin{tikzpicture}[scale=#1,every node/.style={scale=#1}]
  \draw (0,0) rectangle (2.4,2.4);
  \node[anchor=north west] at (0.15,2.25) {Area:};
  \node[anchor=east]  at (-0.15,1.2) {#2};
  \node[anchor=north] at (1.2,-0.15) {#2};
  \begin{scope}[xshift=4.6cm]
    \draw (0,0) rectangle (2.4,2.4);
    \node[anchor=north west] at (0.15,2.25) {Area:};
    \node[anchor=east]  at (-0.15,1.2) {#3};
    \node[anchor=north] at (1.2,-0.15) {#3};
  \end{scope}
\end{tikzpicture}
\end{center}}

\providecommand{\samecubes}[3][1]{%
\begin{center}
{\footnotesize These cubes are the \textbf{same size}.}

\vspace{1.5mm}
\begin{tikzpicture}[scale=#1,every node/.style={scale=#1}]
  \foreach \dx in {0,5.2} {
    \begin{scope}[xshift=\dx cm]
      \draw (0,0) rectangle (2.0,2.0);
      \draw (0,2.0) -- (0.65,2.6) -- (2.65,2.6) -- (2.0,2.0);
      \fill[black!12] (2.0,0) -- (2.65,0.6) -- (2.65,2.6) -- (2.0,2.0) -- cycle;
      \draw (2.0,0) -- (2.65,0.6) -- (2.65,2.6) -- (2.0,2.0);
      \node[anchor=north west] at (0.1,1.9) {Volume:};
    \end{scope}
  }
  \node[anchor=east]  at (-0.15,1.0) {#2};
  \node[anchor=north] at (1.0,-0.15) {#2};
  \node[anchor=west]  at (2.25,0.4) {#2};
  \node[anchor=east]  at (5.05,1.0) {#3};
  \node[anchor=north] at (6.2,-0.15) {#3};
  \node[anchor=west]  at (7.45,0.4) {#3};
\end{tikzpicture}
\end{center}}

%% NOTE: do NOT add \usepackage to individual activity .tex files.
%% When a xourse inlines an activity it gobbles \documentclass and \input, but a bare
%% \usepackage still executes -- after \begin{document} -- and errors out.
%% All shared packages and macros belong here.
%%
%% ximera.cls already loads: amsmath, amssymb, amsthm, cancel, enumitem, graphicx,
%% hyperref, listings, pgfplots, tikz, url, xcolor. No need to re-load those.

\title{Homework 4: Tiered Systems}
\begin{document}
\begin{abstract}
Applying the tax-schedule machinery to two non-tax settings: a discount club's rebate
tiers, and the Pell Grant expected-contribution table. Effective versus marginal rates,
working backwards from a known payout, and reasoning about what a negative tier value
means.
\end{abstract}
\maketitle

\section*{Part 1 --- Discount Club}

Suppose you are part of a Discount Club which provides members with monthly rebates on
purchases made to various companies. To incentivize members to buy, the percent rebate
increases as the total value of purchases increases. The chart below indicates the
rebates the club pays out monthly depending on a member's spending.

\begin{center}
\begin{tabular}{|l|l|l|}
\hline
\textbf{If total purchases are over} & \textbf{but not over} & \textbf{the rebate is} \\
\hline
\$100    & \$800    & $5\%$ of the amount over \$100 \\
\hline
\$800    & \$1,200  & \$35 plus $10\%$ of the amount over \$800 \\
\hline
\$1,200  & \$2,000  & \$75 plus $15\%$ of the amount over \$1,200 \\
\hline
\$2,000  & no limit & \$195 plus $20\%$ of the amount over \$2,000 \\
\hline
\end{tabular}
\end{center}

\begin{problem}
What rebate would a member receive if they spent \$935 in a given month?

\$$\answer[tolerance=0.5]{48.50}$

\begin{freeResponse}
\end{freeResponse}

\begin{solution}
\$935 falls in the second row, so the rebate is \$35 plus $10\%$ of the amount over
\$800:
\[
\$35 + 0.10 \cdot (\$935 - \$800) = \$35 + \$13.50 = \$48.50.
\]
\end{solution}
\end{problem}

\begin{problem}
What would be the effective rebate rate for this spending? Give the calculation, the
mathematical reason for that calculation, and the reason in terms of the context.

About $\answer[tolerance=0.2]{5.19}\%$

\begin{freeResponse}
\end{freeResponse}

\begin{solution}
\textbf{Calculation:}
\[
\frac{\$48.50}{\$935} = 0.05187\ldots \approx 5.19\%.
\]

\textbf{Math reason:} An effective rate compares the total amount received against the
total amount it was computed from --- the whole rebate over the whole spending, not the
rate from any single row.

\textbf{Context reason:} It answers ``across everything I spent this month, what
fraction came back to me?'' Notice it is $5.19\%$, which sits just above the $5\%$ of
the first tier and well below the $10\%$ of the tier this member actually landed in.
That is because most of their spending was rebated at the lower rate, and only the last
\$135 earned $10\%$.
\end{solution}
\end{problem}

\begin{problem}
Suppose a member spent \$1,300 last month and then \$1,600 the next month. What is the
marginal change in total purchases, and what is the marginal change in rebates received?

Change in purchases: \$$\answer{300}$

Rebate at \$1,300: \$$\answer{90}$ \qquad Rebate at \$1,600: \$$\answer{135}$

Change in rebates: \$$\answer{45}$

\begin{freeResponse}
\end{freeResponse}

\begin{solution}
The change in purchases is $\$1{,}600 - \$1{,}300 = \$300$.

To find the change in rebates we must calculate the rebate at each spending level
separately. Both fall in the third row:
\[
\text{rebate}(\$1{,}300) = \$75 + 0.15 \cdot (\$1{,}300 - \$1{,}200) = \$75 + \$15 = \$90
\]
\[
\text{rebate}(\$1{,}600) = \$75 + 0.15 \cdot (\$1{,}600 - \$1{,}200) = \$75 + \$60 = \$135
\]
So the change in rebates is $\$135 - \$90 = \$45$.
\end{solution}
\end{problem}

\begin{problem}
What would be the marginal rebate rate for this change in spending, and why does that
rate make sense given the rebate table?

$\answer{15}\%$

\begin{freeResponse}
\end{freeResponse}

\begin{solution}
\textbf{Calculation:}
\[
\frac{\$45}{\$300} = 0.15 = 15\%.
\]

\textbf{Math reason:} A marginal rate compares the \emph{change} in rebate against the
\emph{change} in spending --- it describes only the additional \$300, not the whole
\$1,600.

\textbf{Context reason:} It answers ``of the extra \$300 I spent, how much came back?''

\textbf{Why it makes sense:} Both \$1,300 and \$1,600 sit inside the same row of the
table, the one covering \$1,200 to \$2,000 at $15\%$. Since the entire \$300 increase
happened within that single tier, the whole increase was rebated at that tier's rate,
and the marginal rate is exactly $15\%$.

Contrast this with the raise that crossed a tax bracket boundary: there the marginal
rate landed between the two bracket rates, because the increase was split across two
tiers. Here nothing was split, so we get a rate straight off the table.
\end{solution}
\end{problem}

\begin{problem}
You overhear another club member bragging that they got \$125 in rebates last month. How
much did they spend? Your explanation should include how and why you know which line of
the rebate table was used.

About \$$\answer[tolerance=15]{1533.33}$

\begin{freeResponse}
\end{freeResponse}

\begin{solution}
First identify the row. The third row spans rebates from \$75 (at \$1,200 of spending)
up to \$195 (at \$2,000). Since $\$75 < \$125 < \$195$, the third row was used. It
cannot be the second row, whose largest rebate is \$75, nor the fourth, whose smallest
is \$195.

Now reverse that row's instruction:
\[
\$125 = \$75 + 0.15 \cdot (P - \$1{,}200)
\]
\[
\$50 = 0.15 \cdot (P - \$1{,}200)
\quad\Longrightarrow\quad
P - \$1{,}200 = \frac{\$50}{0.15} = \$333.33
\]
\[
P = \$1{,}200 + \$333.33 = \$1{,}533.33.
\]
\end{solution}
\end{problem}

\begin{problem}
To be part of the Discount Club there is a monthly fee of \$25. How much do you need to
spend each month for membership to be worth it?

More than \$$\answer[tolerance=10]{600}$

\begin{freeResponse}
\end{freeResponse}

\begin{solution}
Membership is worth it when the rebate exceeds the fee, so we need the spending $P$ at
which the rebate equals \$25 and then go above it.

A \$25 rebate is small, so try the first row, where the rebate is $5\%$ of the amount
over \$100:
\[
\$25 = 0.05 \cdot (P - \$100)
\quad\Longrightarrow\quad
P - \$100 = \$500
\quad\Longrightarrow\quad
P = \$600.
\]
\$600 does fall inside the first row's range of \$100 to \$800, so the assumption was
consistent and the answer stands: you need to spend more than \$600 in a month.

Worth noticing: below \$100 of spending the rebate is \$0, so a member who spends very
little pays \$25 a month for nothing.
\end{solution}
\end{problem}

\section*{Part 2 --- Pell Grants}

Pell Grants are grants, not loans, meaning they do not have to be repaid. They are
awarded to undergraduate students to aid them in paying for college tuition. We assume
we are looking at a Pell Grant for a dependent student. After a questionnaire about
family income, it comes to a table where an amount to be paid by the parents is
determined from the Adjusted Available Income (AAI).

\begin{center}
\begin{tabular}{|l|l|}
\hline
\textbf{If the parents' AAI is} & \textbf{Then the parents' contribution from AAI is} \\
\hline
Less than $-$\$6,820      & $-$\$1,500 \\
\hline
$-$\$6,820 to \$20,600    & $22\%$ of AAI \\
\hline
\$20,601 to \$25,800      & \$4,532 plus $25\%$ of AAI over \$20,600 \\
\hline
\$25,801 to \$31,000      & \$5,832 plus $29\%$ of AAI over \$25,800 \\
\hline
\$31,001 to \$36,300      & \textbf{(A)} plus $34\%$ of AAI over \$31,000 \\
\hline
\$36,301 to \$41,500      & \textbf{(B)} plus $40\%$ of AAI over \$36,300 \\
\hline
\$41,501 or more          & \$11,222 plus $47\%$ of AAI over \$41,500 \\
\hline
\end{tabular}
\end{center}

\subsection*{Borderline AAIs}

\begin{problem}
If a parents' AAI is \$25,800, how much would their contribution be, according to this
table?

\$$\answer{5832}$

\begin{solution}
\$25,800 is the top of the row spanning \$20,601 to \$25,800:
\[
\$4{,}532 + 0.25 \cdot (\$25{,}800 - \$20{,}600) = \$4{,}532 + \$1{,}300 = \$5{,}832.
\]
Notice this is exactly the base amount printed in the \emph{next} row --- the same
continuity pattern as the tax schedule.
\end{solution}
\end{problem}

\begin{problem}
Use what you noticed in the previous question to fill in blanks (A) and (B).

(A) $=$ \$$\answer{7340}$ \qquad (B) $=$ \$$\answer{9142}$

\begin{freeResponse}
\end{freeResponse}

\begin{solution}
Each base is the contribution owed by someone sitting exactly at the bottom edge of that
row, computed from the row below.

\textbf{Blank (A)} is the contribution at an AAI of \$31,000, from the \$25,801--\$31,000
row:
\[
\$5{,}832 + 0.29 \cdot (\$31{,}000 - \$25{,}800) = \$5{,}832 + \$1{,}508 = \$7{,}340.
\]

\textbf{Blank (B)} is the contribution at an AAI of \$36,300, now using the row we have
just completed:
\[
\$7{,}340 + 0.34 \cdot (\$36{,}300 - \$31{,}000) = \$7{,}340 + \$1{,}802 = \$9{,}142.
\]

You can check (B) against the final row: $\$9{,}142 + 0.40 \cdot (\$41{,}500 -
\$36{,}300) = \$9{,}142 + \$2{,}080 = \$11{,}222$, which is the base the table already
prints for the last row.
\end{solution}
\end{problem}

\subsection*{Calculating Contributions}

\begin{problem}
If a parent's AAI is \$28,000, how much are they expected to contribute?

\$$\answer{6470}$

\begin{solution}
\$28,000 falls in the \$25,801--\$31,000 row:
\[
\$5{,}832 + 0.29 \cdot (\$28{,}000 - \$25{,}800) = \$5{,}832 + \$638 = \$6{,}470.
\]
\end{solution}
\end{problem}

\begin{problem}
What is the effective rate of contribution in this scenario? Give the calculation, the
math reason, and the context reason.

About $\answer[tolerance=0.3]{23.11}\%$

\begin{freeResponse}
\end{freeResponse}

\begin{solution}
\textbf{Calculation:}
\[
\frac{\$6{,}470}{\$28{,}000} = 0.23107\ldots \approx 23.11\%.
\]

\textbf{Math reason:} The total contribution divided by the total AAI it was computed
from.

\textbf{Context reason:} It says that across the family's whole adjusted available
income, about $23\%$ of it is expected to go toward college costs. Note this sits below
the $29\%$ marginal rate of the row they are in, for the same reason effective tax
rates sit below marginal ones.
\end{solution}
\end{problem}

\begin{problem}
Is this effective rate analogous to an effective tax rate compared to total income, or
to taxable income? Explain your choice.

\begin{multipleChoice}
\choice{Total income}
\choice[correct]{Taxable income}
\end{multipleChoice}

\begin{freeResponse}
\end{freeResponse}

\begin{solution}
It is analogous to the rate computed against \emph{taxable} income.

The reason is what AAI is: it is the figure that remains after the questionnaire has
already made its adjustments and allowances, and it is the number the table actually
operates on. That is exactly the role taxable income plays in the tax schedule --- the
post-deduction figure the brackets are applied to.

The family's actual gross income, before those adjustments, would be the analogue of
total income, and dividing by that would give a smaller percentage.
\end{solution}
\end{problem}

\begin{problem}
If a parent's contribution is expected to be \$4,000, what was determined to be their
AAI?

About \$$\answer[tolerance=200]{18181.82}$

\begin{freeResponse}
\end{freeResponse}

\begin{solution}
First find the row. The row for AAI from $-$\$6,820 to \$20,600 produces contributions
up to $0.22 \cdot \$20{,}600 = \$4{,}532$ at its top. Since $\$4{,}000 < \$4{,}532$, the
contribution came from this row.

That row is simply $22\%$ of AAI, so
\[
\$4{,}000 = 0.22 \cdot \text{AAI}
\quad\Longrightarrow\quad
\text{AAI} = \frac{\$4{,}000}{0.22} = \$18{,}181.82.
\]
That is within the row's range, confirming the choice.
\end{solution}
\end{problem}

\subsection*{Comparing to Taxes}

\begin{problem}
If a parent's AAI is $-$\$5,000, how much are they expected to contribute? What does
this mean in the ``real world''?

\$$\answer{-1100}$

\begin{freeResponse}
\end{freeResponse}

\begin{solution}
$-$\$5,000 falls in the $-$\$6,820 to \$20,600 row, which is $22\%$ of AAI:
\[
0.22 \cdot (-\$5{,}000) = -\$1{,}100.
\]

A negative expected contribution means the family is not expected to pay anything toward
tuition --- and more than that, the negative figure feeds into the aid formula as
additional demonstrated need, increasing the grant the student is eligible for. An AAI
below zero indicates the family's allowances exceeded their available income.

Note also the floor: the first row caps this at $-$\$1,500 for any AAI below
$-$\$6,820. Checking continuity, $0.22 \cdot (-\$6{,}820) = -\$1{,}500.40$, which is
the $-$\$1,500 of the row above, rounded. So the table stops rewarding an increasingly
negative AAI past that point.
\end{solution}
\end{problem}

\begin{problem}
Describe at least two ways in which this structure is \textbf{similar to} the Federal
Income Tax system as we discussed in class.

\begin{freeResponse}
\end{freeResponse}

\begin{solution}
Several defensible answers; two of the strongest:

\emph{The base-plus-rate structure is identical.} Each row is ``some fixed amount plus a
percentage of the amount over the row's lower boundary,'' and each fixed amount is the
total accumulated from all rows below --- which is why we could compute blanks (A) and
(B) by exactly the method we used for blanks A, B, and C in the tax schedule.

\emph{It is continuous, so it has marginal rates rather than cliffs.} Crossing a boundary
changes the rate only on the amount above it. A family whose AAI rises by a dollar never
sees their expected contribution jump.

\emph{The calculation is reversible.} Given a contribution, you can identify the row from
the range of contributions it can produce, then undo the arithmetic to recover the AAI ---
exactly as we did working backwards from a tax bill.
\end{solution}
\end{problem}

\begin{problem}
Describe at least two ways in which this structure is \textbf{different from} the
Federal Income Tax system as we discussed in class.

\begin{freeResponse}
\end{freeResponse}

\begin{solution}
Again several defensible answers:

\emph{Values can be negative.} There is no such thing as negative taxable income in the
tax schedule, but AAI can be negative, and the table explicitly handles it --- and even
places a floor of $-$\$1,500 on the result. The tax schedule has no analogous floor.

\emph{The rates climb much faster over a much narrower range.} The tax schedule takes
until \$510,300 to reach $37\%$. This table reaches $47\%$ by \$41,501. The tiers are
only about \$5,200 wide, compared to tax brackets tens of thousands of dollars wide.

\emph{The direction of benefit is reversed.} A higher tax means money leaving you for the
government; a higher expected contribution means less grant aid. Both are computed the
same way, but one is a payment and the other is a reduction in a benefit.

\emph{The boundaries are stated in whole dollars} (\$20,600 then \$20,601) rather than
the tax schedule's ``over\ldots but not over'' phrasing, which technically leaves
fractional amounts between rows undefined.
\end{solution}
\end{problem}

\end{document}
